\documentclass[11pt, a4paper, oneside]{article}
\usepackage[ngerman]{babel}
\usepackage{worksheet}

\begin{document}
	\author{L. Bung}
	\title{Exponentialfunktionen: \hspace{10cm} Einführung}
	\subject{Mathematik}
	\class{FHR1}
	\maketitle
	
	\partnertask{Papier falten}
	
	Nehmen Sie ein DIN-A4-Papier zur Hand.
	
	a) Wie häufig können Sie das Papier in der Mitte falten?
	
	\checkered[1.5cm]
	
	b) Aus wie vielen Schichten besteht Ihr Papier am Ende?
	
	\checkered[3cm]
	
	c) Wie viele Schichten wären es, wenn man 10 mal / 15 mal / $n$ mal falten würde?
	
	\checkered[3cm]
	
	d) Ein durchschnittliches Papier ist etwa 0,1 mm dick.
	Wie oft muss man es falten, um einen Stapel bis zum Mond (Entfernung: etwa 384 000 km) zu erreichen?
	
	\checkered[3cm]
	
	\begin{hint}{Lineares und exponentielles Wachstum}
		\phantom{x}\vspace{6cm}
	\end{hint}
	
	\singletask{Linear oder exponentiell?}
	
	Entscheiden Sie, ob es sich um lineares oder exponentielles Wachstum handelt.
	Begründen Sie!
	
	\begin{multicols}{3}
		a)
		\begin{tabular}{|c||c|c|c|c|}
			\hline
			$x$ & 0 & 1 & 2 & 3\\
			\hline
			$f(x)$ & 8 & 12 & 18 & 27\\
			\hline
		\end{tabular}
		
		b)
		\begin{tabular}{|c||c|c|c|c|}
			\hline
			$x$ & 0 & 1 & 2 & 3\\
			\hline
			$f(x)$ & 7 & 11 & 15 & 19\\
			\hline
		\end{tabular}
		
		c)
		\begin{tabular}{|c||c|c|c|c|}
			\hline
			$x$ & 0 & 1 & 2 & 3\\
			\hline
			$f(x)$ & 32 & 16 & 8 & 4\\
			\hline
		\end{tabular}
	\end{multicols}
	
	\checkered[6cm]
	
	d) Füllen Sie die fehlenden Werte so aus, dass es sich um exponentielles Wachstum handelt.
	
	\begin{table}[ht]
		\centering
		\begin{tabularx}{.5\textwidth}{|c||X|X|X|X|X|X|}
			\hline
			$x$ & 0 & 1 & 2 & 3 & 4 & 5\\
			\hline
			$f(x)$ &&& 2,5 & 3,5 &&\\
			\hline
		\end{tabularx}
	\end{table}
	
	
%	\singletask{Industriestanzer}
%	
%	Eine Firma stellt Maschinen her, welche zum industriellen Stanzen von Verpackungsmaterial verwendet werden.
%	
%	a) Eine Maschine eignet sich zum Stanzen von 35 Lagen Papier.
%	Wie häufig dürfte man ein Papier falten, damit es gerade noch gestanzt werden kann?
%	
%	\checkered[2cm]
%	
%	b) Wie viele Lagen müsste die Maschine mindestens stanzen können, damit man das Papier 6 mal falten kann?
%	
%	\checkered[1.5cm]
%	
%	c) Faltet man das Papier wie einen Brief, sind es jeweils drei Lagen.
%	Wie häufig darf man das Papier so falten, damit die Maschine es noch stanzen kann?
%	
%	\checkered[1.5cm]
	
	\grouptask{Die Exponentialfunktion und ihre Parameter}
	
	\begin{wrapfigure}{r}{.25\textwidth}
		\vspace{-\baselineskip}
		\qrlink{https://www.geogebra.org/m/sucjkxse}
	\end{wrapfigure}
	
	Eine Exponentialfunktion hat die allgemeine Form $f(x) = a \cdot q^{b(x-c)} + d$.
	Finden Sie mithilfe der Geogebra-App heraus, welche Bedeutung die Parameter $q$, $a$, $b$, $c$ und $d$ haben.
	
	Füllen Sie dazu die folgende Tabelle aus.
	
	\vspace{1cm}
	
	\begin{table}[!ht]
		\begin{tabularx}{\textwidth}{|l|X|}
			\hline
			\textbf{Parameter} & \textbf{Beschreibung}\\
			\hline
			\hline
			$q$&\\[1.5cm]
			\hline
			$a$&\\[1.5cm]
			\hline
			$b$&\\[1.5cm]
			\hline
			$c$&\\[1.5cm]
			\hline
			$d$&\\[1.5cm]
			\hline
		\end{tabularx}
	\end{table}
	
	\begin{goallist}
		...Unterschiede zwischen linearem und exponentiellem Wachstum beschreiben.&&&\\
		\hline
		...entscheiden, ob Wachstum linear oder exponentiell ist.&&&\\
		\hline
		...die Bedeutung der Basis $q$ für das Wachstum angeben.&&&\\
		\hline
		...den Einfluss der Parameter auf den Graphen der Exponentialfunktion erklären.&&&\\
	\end{goallist}
\end{document}